\documentclass[10pt,a4paper]{amsart}
\usepackage[margin=19mm]{geometry}
\usepackage{amsmath,amssymb,mathtools}
\usepackage[scr=rsfs,cal=euler]{mathalfa}
\usepackage{xfrac}
\usepackage[unicode,hidelinks,bookmarks=false]{hyperref}
\newcommand{\ie}{i.e.\ }
\newcommand{\cf}{cf.\ }
\newcommand{\resp}{respectively\ }
\newcommand{\Subsection}{\subsection}
\providecommand{\nobreakdash}{\nobreak\hbox{-}}
\setlength{\emergencystretch}{2em}
\multlinegap0pt
\usepackage{bm}
\usepackage{bbm}
\usepackage{dsfont}
\usepackage{xspace}
\usepackage{cancel}
\usepackage{mathtools}
\usepackage{amsthm}
\makeatletter
\@ifundefined{assumption}{\newtheorem{assumption}{Assumption}}{}
\@ifundefined{theorem}{\newtheorem{theorem}{Theorem}}{}
\makeatother
\def\H{\mathcal{H}}
\newcommand{\dd}{\,\mathrm{d}}
\newcommand{\inner}[2]{\langle #1, #2 \rangle}
\title[Polynomial Eigenfunctions and Matrix Lyapunov Equations from]{Polynomial Eigenfunctions and Matrix Lyapunov Equations from Energy Balance Integrals}
\author{Netzer Moriya}
\date{Source: arXiv:2506.15288v1 (CC BY 4.0)}
\begin{document}
\maketitle

\noindent\emph{Source and licence.} Polynomial Eigenfunctions and Matrix Lyapunov Equations from Energy Balance Integrals. Version arXiv:2506.15288v1 (CC BY 4.0), distributed under the Creative Commons Attribution 4.0 International licence, as recorded in the arXiv licence metadata for this exact version and shown on the abstract page https://arxiv.org/abs/2506.15288v1. Full rights notice: licenses/arxiv-2506.15288-CC-BY-4.0.txt.\par\smallskip

\noindent\textit{Editorial scope.} Counted unit: the Theorem whose source label is \texttt{thm:covariance\_properties} in the article (source0.tex lines 215--223, source label \texttt{thm:covariance\_properties}), delivered together with the article's own complete original proof of it (source0.tex lines 225--258, one \texttt{proof} environment of 34 lines). No mathematics is written, repaired or re-derived: statement and proof are the article's own text line for line. Required context retained uncounted: ass:selfadjoint (lines 128--139); eq:energy\_integral (lines 155--157); ass:semigroup (lines 173--179); thm:rigorous\_convergence (lines 198--203). Every definition, display and label of those lines is retained unchanged. Omitted from this package and not redistributed: 1--127, 140--154, 158--172, 180--197, 204--214, 224--224, 259--728. Nothing inside a retained excerpt is omitted or altered: every retained body is delivered exactly as the article's lines are written, so no omission receipt of any kind accompanies this package. Citations: the retained excerpts carry no citation command at all, so no bibliography entry of the article is transcribed and no reference is redistributed. Reference adaptation: none was needed and none was made; every reference inside the delivered unit resolves inside it, so no unresolved reference or citation is produced. Printed slips kept literal: no printed word, symbol, hypothesis or number of the counted statement or of its proof was altered. Layout and class: the article's own class is \texttt{article} with packages including amsmath, amsfonts, amsthm, amssymb, inputenc, fontenc, graphicx, caption, color, hyperref, breakurl, grffile, makeidx, acronym; the wrapper is the lane's pinned \texttt{amsart} editorial wrapper at 10pt on A4, which restates the theorem-like environments the retained text needs and adds the article's own macro definitions that the retained text needs (\texttt{\textbackslash H}, \texttt{\textbackslash dd}, \texttt{\textbackslash inner}), with extra wrapper packages (bm, bbm, dsfont, xspace, cancel, mathtools). The delivered numbering is the unit's own. Assets: no retained span contains a figure environment, a graphics-inclusion command, a PDF-page inclusion, a table or a TikZ picture; no image file is copied into this package, the package declares no asset dependency and no image file is redistributed with it. Licence: version arXiv:2506.15288v1 of this article is distributed under the Creative Commons Attribution 4.0 International licence, as recorded in the arXiv licence metadata for this exact version and shown on the abstract page https://arxiv.org/abs/2506.15288v1. A full rights notice is delivered at licenses/arxiv-2506.15288-CC-BY-4.0.txt. Characters: every retained excerpt is pure ASCII, so the delivered engine, XeLaTeX with its default Latin Modern text font, reports no missing character.
\begin{assumption}[Self-adjoint dissipative structure]\label{ass:selfadjoint}
The operator $\mathcal{L}$ is self-adjoint and dissipative on $\H$:
\begin{enumerate}
\item \textbf{Self-adjointness}: $\mathcal{L} = \mathcal{L}^*$
\item \textbf{Dissipativity}: $\inner{\phi}{\mathcal{L}\phi} \leq -\gamma \|\phi\|^2$ for some $\gamma > 0$ and all $\phi \in \mathcal{D}(\mathcal{L})$
\item \textbf{Spectral completeness}: $\mathcal{L}$ has a complete orthonormal system of eigenfunctions $\{\phi_k\}_{k=1}^{\infty}$ with real eigenvalues $\{\lambda_k\}_{k=1}^{\infty}$:
\begin{equation}
\mathcal{L}\phi_k = \lambda_k \phi_k, \quad \inner{\phi_i}{\phi_j} = \delta_{ij}, \quad \lambda_k < 0
\end{equation}
\end{enumerate}
\textbf{Consequence of self-adjointness}: Since $\mathcal{L} = \mathcal{L}^*$, all eigenvalues are real ($\lambda_k \in \mathbb{R}$), and $\mathcal{L}^*\phi_k = \mathcal{L}\phi_k = \lambda_k \phi_k$. Therefore, $\lambda_k^* = \lambda_k$ throughout this work.
\end{assumption}

\begin{equation}\label{eq:energy_integral}
\inner{\phi}{P\psi} = \int_0^{\infty} \inner{\phi}{e^{\mathcal{L}t} Q e^{\mathcal{L}t} \psi} \dd t
\end{equation}

\begin{assumption}[Semigroup generation]\label{ass:semigroup}
The operator $\mathcal{L}$ satisfying Assumption \ref{ass:selfadjoint} generates a strongly continuous semigroup $\{e^{\mathcal{L}t}\}_{t \geq 0}$ of contractions on $\H$:
\begin{equation}
\|e^{\mathcal{L}t}\| \leq e^{-\gamma t} \quad \text{for all } t \geq 0
\end{equation}
where $\gamma > 0$ is the dissipation constant from Assumption \ref{ass:selfadjoint}.
\end{assumption}

\begin{theorem}[Rigorous convergence of energy integral]\label{thm:rigorous_convergence}
Under Assumptions \ref{ass:selfadjoint}-\ref{ass:semigroup}, the energy dissipation integral converges absolutely:
\begin{equation}
\left|\int_0^{\infty} \inner{\phi}{e^{\mathcal{L}t} Q e^{\mathcal{L}t} \psi} \dd t\right| \leq \frac{\|\phi\|\|\psi\|\|Q\|}{2\gamma}
\end{equation}
\end{theorem}

\begin{theorem}[Fundamental properties of the energy covariance]\label{thm:covariance_properties}
The energy covariance operator $P$ satisfies:
\begin{enumerate}
\item \textbf{Positive definiteness}: $\inner{\phi}{P\phi} \geq 0$ for all $\phi \in \H$
\item \textbf{Self-adjointness}: $P^* = P$
\item \textbf{Energy balance equation}: $\mathcal{L}P + P\mathcal{L} = -Q$
\item \textbf{Convergence}: The integral \eqref{eq:energy_integral} converges absolutely
\end{enumerate}
\end{theorem}

\begin{proof}
\textit{Positive definiteness}: For any $\phi \in \H$,
\begin{align}
\inner{\phi}{P\phi} &= \int_0^{\infty} \inner{\phi}{e^{\mathcal{L}t} Q e^{\mathcal{L}t} \phi} \dd t\\
&= \int_0^{\infty} \inner{e^{\mathcal{L}t}\phi}{Q e^{\mathcal{L}t}\phi} \dd t \geq 0
\end{align}
since $Q$ is positive.

\textit{Self-adjointness}: Direct from the symmetry of the integral representation and the fact that $Q^* = Q$.

\textit{Energy balance}: Starting with the definition:
\begin{equation}
\inner{\phi}{P\psi} = \int_0^{\infty} \inner{\phi}{e^{\mathcal{L}\tau} Q e^{\mathcal{L}\tau} \psi} \dd\tau
\end{equation}
we compute:
\begin{align}
\inner{\phi}{(\mathcal{L}P + P\mathcal{L})\psi} &= \inner{\mathcal{L}\phi}{P\psi} + \inner{\phi}{P\mathcal{L}\psi}\\
&= \int_0^{\infty} \left(\inner{\mathcal{L}\phi}{e^{\mathcal{L}\tau} Q e^{\mathcal{L}\tau} \psi} + \inner{\phi}{e^{\mathcal{L}\tau} Q e^{\mathcal{L}\tau} \mathcal{L}\psi}\right) \dd\tau\\
&= \int_0^{\infty} \inner{\phi}{(\mathcal{L}e^{\mathcal{L}\tau} Q e^{\mathcal{L}\tau} + e^{\mathcal{L}\tau} Q e^{\mathcal{L}\tau}\mathcal{L})\psi} \dd\tau
\end{align}
where we used self-adjointness: $\inner{\mathcal{L}\phi}{\chi} = \inner{\phi}{\mathcal{L}\chi}$. Since:
\begin{equation}
\frac{\dd}{\dd\tau}(e^{\mathcal{L}\tau} Q e^{\mathcal{L}\tau}) = \mathcal{L}e^{\mathcal{L}\tau} Q e^{\mathcal{L}\tau} + e^{\mathcal{L}\tau} Q e^{\mathcal{L}\tau}\mathcal{L}
\end{equation}
we obtain:
\begin{align}
\inner{\phi}{(\mathcal{L}P + P\mathcal{L})\psi} &= \int_0^{\infty} \inner{\phi}{\frac{\dd}{\dd\tau}(e^{\mathcal{L}\tau} Q e^{\mathcal{L}\tau})\psi} \dd\tau\\
&= \inner{\phi}{\left[e^{\mathcal{L}\tau} Q e^{\mathcal{L}\tau}\right]_0^{\infty}\psi}\\
&= \lim_{\tau \to \infty} \inner{\phi}{e^{\mathcal{L}\tau} Q e^{\mathcal{L}\tau}\psi} - \inner{\phi}{Q\psi}
\end{align}
Since $\lambda_k < 0$ for all $k$, the limit vanishes, giving $\mathcal{L}P + P\mathcal{L} = -Q$.

\textit{Convergence}: Follows from Theorem \ref{thm:rigorous_convergence}.
\end{proof}

\end{document}
